% !TeX root = ../main.tex
% !TeX program = xelatex
% !TeX encoding = UTF-8
\section{多目标规划与物理反馈求解}

\subsection{费用与系统碳的容差端点及约束扫描}

技术、管网和运行模式是离散选择，能流、容量与热状态为连续变量，共同构成混合整数线性规划（MILP）。费用和系统碳排放由同一份运行量计算，目标为
\begin{equation}
F_C=C_0+C_{\mathrm{carbon}}+C_{\mathrm{pref}},\qquad F_E=\boldsymbol a_E^\mathsf T\boldsymbol x+E_0,
\label{eq:scoring-objectives}
\end{equation}
其中$C_0$为基本费用，$C_{\mathrm{carbon}}$为碳费用，$C_{\mathrm{pref}}$为偏好与惩罚项；$\boldsymbol x$为决策变量，$\boldsymbol a_E$为排放系数，$E_0$为排放常数项。各方案共用可靠性、舒适、设备和资源约束，不能靠少供热降低排放。

费用相同仍可能有不同排放。先优化费用，再容许小幅变化以改善排放，得到容差端点；交换目标次序并新建模型求另一端点\supcite{5}。以$f_1,f_2$表示主、次目标，第一阶段得到主目标可行值$v_1$后，第二阶段求解
\begin{equation}
f_1\leq v_1+\delta_{\mathrm{abs}}+\delta_{\mathrm{rel}}|v_1|,
\qquad\min f_2.
\label{eq:scoring-endpoint}
\end{equation}
$\delta_{\mathrm{abs}}$与$\delta_{\mathrm{rel}}$为绝对和相对容差，默认值为$10^{-7}$与$10^{-8}$，前者量纲随主目标。两阶段分别保存状态、时限、目标上下界及相对差距gap。若第一阶段只有限时可行解，第二阶段只说明在该值附近改善次目标，不能声称最优端点。

随后在端点排放区间内设置多档碳帽，即排放上限$\varepsilon_E$，逐档求解$\min F_C$、$F_E\leq\varepsilon_E^{\mathrm{eff}}$，其中$\varepsilon_E^{\mathrm{eff}}=\varepsilon_E+10^{-7}+10^{-8}|\varepsilon_E|$。专用扫描将碳价置零，赛事入口保留配置碳价与偏好项，结果分别标注。各档检查碳帽和数值可行性，去重后剔除支配点：若另一方案两项目标均不更高、至少一项更低，该点即被支配。所得为有限采样非支配集合。

\subsection{实际日选取与硬能量约束下的权重校正}

代表日近似较长时序，权重表示该日代表多少天。将同源气象、负荷按日展开、标准化，用K-means分组，取各组距中心最近的实际日。另保留系统及各服务负荷峰值、最大真实相邻小时爬坡、最低气温和已配置的资源风险日；跨午夜最大爬坡保留两侧日期，新增极值日频数从原组扣除，保持总日数。

频数权重的总能量与负荷持续分布可能偏离原序列，故固定日期，只校正权重。设保留观测有$D$个完整日，每日$H=24$小时；$q_{jh}$为代表日$j$第$h$小时负荷，$E_j=\sum_hq_{jh}\Delta t$为当日能量。两组负荷值合并、去重、排序得到支撑点$s_i$；$F_o(s_i)$和$F_w(s_i)$分别为观测与代表负荷不超过$s_i$的小时比例。求解
\begin{equation}
\begin{aligned}
\min_{\boldsymbol w}\;&\frac{1}{\bar q}\sum_i(s_{i+1}-s_i)
       |F_w(s_i)-F_o(s_i)|,\\
\text{s.t.}\;&\sum_jw_j=D,\quad \sum_jw_jE_j=E_o,\quad w_j\geq\ell_j,\\
F_w(s_i)=\;&\frac{1}{HD}\sum_jw_j\sum_{h=1}^{H}\mathbf1(q_{jh}\leq s_i).
\end{aligned}
\label{eq:scoring-hard-w1}
\end{equation}
其中$w_j$为待求权重，$\ell_j$为下限，$E_o$与$\bar q>0$为保留序列总能量和平均负荷，$\mathbf1$为条件成立时取1的指示函数。目标为一阶Wasserstein距离（$W_1$）除以平均负荷，累计分布曲线间差距，越小表示负荷持续分布越接近。日数与能量等式保持天数和供能总量。绝对值用辅助变量线性化，形成HiGHS求解的线性规划（LP）；再独立检查日数、能量、下限和目标残差。六楼极值日取$\ell_j=1$日，其余取$10^{-6}$日，不改观测与日轮廓。

结果须通过预设检查：相对能量误差不超过0.1\%，同口径归一$W_1$不劣于频数基线，保留序列峰值、真实相邻爬坡及原始观测峰值保留率均为1。失败时保存原因并回退频数权重；回退结果仍须单独检查，失败保留未通过状态。统计压缩通过后仍须连续回放检验热状态。

\subsection{固定投资原时序回放与管网反馈}

代表日权重不能还原跨日室温和储热衔接。回放固定规划选定的技术、容量、管网、储热及接入方案，只重求逐时调度。核对输入校验值、时间戳、容量与服务要求后，连续递推热状态，独立复算舒适、能量和初末条件。三城共同2880小时窗口及控制策略试验分别记录，仅检验所选资产在相应窗口的运行可行性。

管网反馈检查小时平均调度量对应的瞬时输送条件，限定独立泵星型支路、物性固定的水和有限个固定流量、回水温度的稳态模式。各模式计算供回水散热、远端供水温度、摩阻及泵耗。令$z_{mt}$为模式$m$在时段$t$是否启用的二元变量，$q_{mt}$为平均入口热功率，则
\begin{equation}
0\leq q_{mt}\leq Q_m^{\mathrm{in}}z_{mt},\quad
\sum_{m\in b}z_{mt}\leq1;
\qquad z_{mt}=0\quad(m\in\mathcal B^{(k)},\ \forall t),
\label{eq:scoring-network-cut}
\end{equation}
其中$Q_m^{\mathrm{in}}$为模式入口热功率，$b$为支路，$\mathcal B^{(k)}$为第$k$轮不合格模式。交付量按小时运行比例折算；温度、流速、压降、泵功率及源网容量按同时开启时校核，避免平均值掩盖瞬时超限。各轮重算热量和泵耗，全时段排除不合格模式；无新增排除、求解失败或轮数耗尽时停止并保存记录。有效反馈至少排除一个模式；温差软惩罚仅用于偏好排序，物理达标依靠硬校核与排除。模型未计输送时滞、启停瞬态和共泵水力平衡。

\subsection{目标收敛的检验与各步骤衔接}

核心模型经SciPy的\path{milp}接口调用HiGHS\supcite{6,7}，保存终止状态、可行解、矩阵及整数残差、重算目标值和gap。对最小化问题，已找到的可行目标值$U$给出最优值的上界，求解器下界$L$表示最优值不会更低；gap是求解器报告的相对差距，越小表示已知方案越接近最优。二者须对应同一模型、阶段和目标，并同时加回常数项，才能区分限时可行解与容差内最优解。物理校核另行记录，因为目标收敛并不代替运行条件检查。

固定设备热泵LP将计算机以双精度浮点（binary64）格式保存的系数视为精确有理数，利用弱对偶构造电量费下界：可行对偶解的目标值不会高于最优值。配套数值解按给定容差检查可行性。该证书只对应固定资产的编码模型及电量费目标，用于同资产调度比较。
